\[
\textbf{Solution.}
\]

Assume throughout that the negative powers appearing in $Q$ are meaningful, i.e. that $\lambda_1,\lambda_2,\lambda_3\neq 0$. Since $\lambda_3=-\lambda_1-\lambda_2$, this means in particular that $\lambda_1,\lambda_2,\lambda_1+\lambda_2$ are nonzero.

We are given

\[
l_1+l_2+l_3=-2,\qquad l_1\ge l_2\ge l_3,\qquad l_1\le 0.
\]

Because $l_1\ge l_2\ge l_3$, we have

\[
l_2\le l_1\le 0,\qquad l_3\le l_2\le 0.
\]

Thus all three integers $l_1,l_2,l_3$ are nonpositive. If $l_1\le -1$, then

\[
l_1\le -1,\qquad l_2\le -1,\qquad l_3\le -1,
\]

so

\[
l_1+l_2+l_3\le -3,
\]

contradicting $l_1+l_2+l_3=-2$. Therefore

\[
l_1=0.
\]

Then

\[
l_2+l_3=-2,
\]

with

\[
0=l_1\ge l_2\ge l_3.
\]

Since $l_3=-2-l_2$, the condition $l_2\ge l_3$ gives

\[
l_2\ge -2-l_2,
\]

hence

\[
2l_2\ge -2,
\]

so

\[
l_2\ge -1.
\]

Also $l_2\le 0$. Since $l_2\in\mathbb Z$, the only possibilities are

\[
l_2=0 \quad\text{or}\quad l_2=-1.
\]

If $l_2=0$, then $l_3=-2$, giving

\[
(l_1,l_2,l_3)=(0,0,-2).
\]

If $l_2=-1$, then $l_3=-1$, giving

\[
(l_1,l_2,l_3)=(0,-1,-1).
\]

Thus the only possible triples are

\[
(l_1,l_2,l_3)\in\{(0,0,-2),(0,-1,-1)\}.
\]

Now we verify the two cases.

---

### Case $1$: \((l_1,l_2,l_3)=(0,0,-2)\)

The parity-restricted sums are

\[
\sum_{\substack{1\le k\le l_1\\ k\equiv l_1\pmod 2}} A(\ldots,k)
=
\sum_{\substack{1\le k\le 0\\ k\equiv 0\pmod 2}} A(\ldots,k),
\]

and

\[
\sum_{\substack{1\le k\le l_2\\ k\equiv l_2\pmod 2}} B(\ldots,k)
=
\sum_{\substack{1\le k\le 0\\ k\equiv 0\pmod 2}} B(\ldots,k).
\]

Both indexing sets are empty, so both sums are equal to $0$.

Next compute $Q$. We have

\[
\bar l_1=\bar 0=0,\qquad \bar l_2=\bar 0=0,
\]

and since $l_3=-2<0$,

\[
\bar l_3=-l_3-1=2-1=1.
\]

Also

\[
\sigma(0)=0,\qquad \sigma(1)=\frac{1\cdot 2}{2}=1.
\]

Moreover,

\[
l_1l_2=0\cdot 0=0,\qquad
l_2l_3=0\cdot (-2)=0,\qquad
l_1l_3=0\cdot (-2)=0.
\]

Therefore

\[
Q
=
0\cdot \lambda_1^{-2}
+
0\cdot \lambda_2^{-2}
+
1\cdot \lambda_3^{-2}
+
0
+
0
+
0,
\]

so

\[
Q=\lambda_3^{-2}.
\]

Hence

\[
\lambda_1\lambda_2\lambda_3 Q
=
\lambda_1\lambda_2\lambda_3\cdot \lambda_3^{-2}
=
\lambda_1\lambda_2\lambda_3^{-1}.
\]

On the other hand,

\[
\lambda_1^{l_1+1}\lambda_2^{l_2+1}\lambda_3^{l_3+1}
=
\lambda_1^{0+1}\lambda_2^{0+1}\lambda_3^{-2+1}
=
\lambda_1\lambda_2\lambda_3^{-1}.
\]

Thus

\[
\lambda_1^{l_1+1}\lambda_2^{l_2+1}\lambda_3^{l_3+1}
=
\lambda_1\lambda_2\lambda_3 Q.
\]

---

### Case $2$: \((l_1,l_2,l_3)=(0,-1,-1)\)

The first parity-restricted sum is

\[
\sum_{\substack{1\le k\le l_1\\ k\equiv l_1\pmod 2}} A(\ldots,k)
=
\sum_{\substack{1\le k\le 0\\ k\equiv 0\pmod 2}} A(\ldots,k),
\]

which is empty. The second is

\[
\sum_{\substack{1\le k\le l_2\\ k\equiv l_2\pmod 2}} B(\ldots,k)
=
\sum_{\substack{1\le k\le -1\\ k\equiv -1\pmod 2}} B(\ldots,k),
\]

which is also empty. Hence both sums are $0$.

Now compute $Q$. We have

\[
\bar l_1=\bar 0=0.
\]

Since $l_2=l_3=-1<0$,

\[
\bar l_2=-l_2-1=1-1=0,
\]

and

\[
\bar l_3=-l_3-1=1-1=0.
\]

Therefore

\[
\sigma(\bar l_1)=\sigma(\bar l_2)=\sigma(\bar l_3)=\sigma(0)=0.
\]

The pairwise products are

\[
l_1l_2=0\cdot (-1)=0,
\]

\[
l_2l_3=(-1)(-1)=1,
\]

and

\[
l_1l_3=0\cdot (-1)=0.
\]

Thus

\[
Q
=
0
+
0
+
0
+
0
+
1\cdot \lambda_2^{-1}\lambda_3^{-1}
+
0,
\]

so

\[
Q=\lambda_2^{-1}\lambda_3^{-1}.
\]

Hence

\[
\lambda_1\lambda_2\lambda_3 Q
=
\lambda_1\lambda_2\lambda_3\cdot \lambda_2^{-1}\lambda_3^{-1}
=
\lambda_1.
\]

On the other hand,

\[
\lambda_1^{l_1+1}\lambda_2^{l_2+1}\lambda_3^{l_3+1}
=
\lambda_1^{0+1}\lambda_2^{-1+1}\lambda_3^{-1+1}
=
\lambda_1\lambda_2^0\lambda_3^0
=
\lambda_1.
\]

Therefore

\[
\lambda_1^{l_1+1}\lambda_2^{l_2+1}\lambda_3^{l_3+1}
=
\lambda_1\lambda_2\lambda_3 Q.
\]

---

Combining the enumeration with the two direct verifications, we conclude that if

\[
l_1+l_2+l_3=-2,\qquad l_1\ge l_2\ge l_3,\qquad l_1\le 0,
\]

then necessarily

\[
(l_1,l_2,l_3)\in\{(0,0,-2),(0,-1,-1)\},
\]

both parity-restricted sums defining $R$ are empty, so

\[
R=0,
\]

and in both cases

\[
\lambda_1^{l_1+1}\lambda_2^{l_2+1}\lambda_3^{l_3+1}
=
\lambda_1\lambda_2\lambda_3 Q.
\]

This proves the statement.